\documentclass[../../../include/open-logic-section]{subfiles}

\begin{document}

\olfileid{sth}{cardinals}{classing}
\olsection{متناهی، \usetoken{S}{enumerable}، \usetoken{S}{nonenumerable}}

اکنون که با کاردینال‌ها آشنا شده‌ایم، شایسته است اندکی دربارهٔ گونه‌های
مختلف آن‌ها، به‌ویژه کاردینال‌های متناهی، !!{enumerable} و
!!{nonenumerable}، سخن بگوییم.

دو نتیجهٔ نخست ما مستلزم آن‌اند که کاردینال‌های متناهی دقیقاً همان اعداد
ترتیبی متناهی باشند که پیش‌تر در~\olref[z][infinity-again]{defnomega}
آن‌ها را \emph{اعداد طبیعی} تعریف کردیم:

\begin{prop}\ollabel{finitecardisoequal}
فرض کنید~$n,m\in\omega$. آنگاه~$n=m$ اگر و تنها اگر
$\cardeq{n}{m}$.
\end{prop}

\begin{proof}
جهت \emph{چپ به راست} بدیهی است. برای اثبات جهت \emph{راست به چپ}،
فرض کنید با وجود~$n\neq m$، $\cardeq{n}{m}$. بنا بر سه‌شقی، یا
$n\in m$ یا~$m\in n$؛ بدون کاستن از کلیت فرض کنید~$n\in m$. آنگاه
$n\subsetneq m$ و یک~!!a{bijection} به‌صورت~$f\colon m\to n$ وجود
دارد؛ پس~$m$ نامتناهیِ ددکیندی است، برخلاف
\olref[z][infinity-again]{naturalnumbersarentinfinite}.
\end{proof}

\begin{cor}\ollabel{naturalsarecardinals}
اگر~$n\in\omega$، آنگاه~$n$ یک کاردینال است.
\end{cor}

\begin{proof}
بی‌درنگ.
\end{proof}
\noindent
همچنین نتیجه می‌شود چند برداشت معقول از معنای «متناهی» یا «نامتناهی»
بودن یک کاردینال بر هم منطبق‌اند:
\begin{thm}\ollabel{generalinfinitycharacter}
برای هر مجموعهٔ~$A$، عبارت‌های زیر هم‌ارزند:
\begin{enumerate}
	\item\ollabel{card:notinomega} $\card{A}\notin\omega$؛ یعنی
	$\card{A}$ یک عدد طبیعی نیست؛
	\item\ollabel{card:omegaplus} $\omega\leq\card{A}$؛
	\item\ollabel{card:infinite} $A$ نامتناهیِ ددکیندی است.
\end{enumerate}
\end{thm}

\begin{proof}
از~\olref[ord-arithmetic][using-addition]{ordinfinitycharacter}،
\olref[cardinals][cardsasords]{lem:CardinalsBehaveRight} و
\olref{naturalsarecardinals}.
\end{proof}

این نتیجه تعریف \emph{صوریِ} مفاهیمی را مجاز می‌سازد که در
\olref[sfr][][]{part} نسبتاً غیرصوری به کار بردیم:

\begin{defn}\ollabel{defnfinite}
می‌گوییم~$A$ \emph{متناهی} است اگر و تنها اگر~$\card{A}$ یک عدد طبیعی
باشد؛ یعنی~$\card{A}\in\omega$. در غیر این صورت می‌گوییم~$A$
\emph{نامتناهی} است.
\end{defn}
\noindent
اما توجه کنید که این تعریف با فرض زمینه‌ایِ~$\ZFC$ عرضه می‌شود. زیرا
برای تضمین اینکه هر مجموعه کاردینالیته‌ای دارد، به اصل خوش‌ترتیبی نیاز
داشتیم. در واقع، بدون خوش‌ترتیبی ممکن است مجموعه‌ای نه متناهی و نه
نامتناهیِ ددکیندی باشد. در~\olref[choice][]{chap} به این نوع مسائل
بازخواهیم گشت. فعلاً همچنان بر خوش‌ترتیبی تکیه می‌کنیم.

اکنون از کاردینال‌های متناهی به کاردینال‌های نامتناهی بپردازیم. نخست دو
نکتهٔ مقدماتی:

\begin{cor}\ollabel{omegaisacardinal}
$\omega$ کوچک‌ترین کاردینال نامتناهی است.
\end{cor}

\begin{proof}
$\omega$ یک کاردینال است، زیرا نامتناهیِ ددکیندی است؛ و اگر برای یک
$n\in\omega$، $\cardeq{\omega}{n}$، آنگاه~$n$ نیز نامتناهیِ ددکیندی
می‌بود، برخلاف~\olref[z][infinity-again]{naturalnumbersarentinfinite}.
اکنون بنا بر تعریف، $\omega$ کوچک‌ترین کاردینال نامتناهی است.
\end{proof}

\begin{cor}
هر کاردینال نامتناهی یک عدد ترتیبی حدی است.
\end{cor}

\begin{proof}
فرض کنید~$\alpha$ یک عدد ترتیبی جانشینِ نامتناهی باشد؛ پس برای یک
$\beta$، $\alpha=\beta\ordplus1$. بنا بر~\olref{finitecardisoequal}،
$\beta$ نیز نامتناهی است؛ ازاین‌رو بنا بر
\olref[ord-arithmetic][using-addition]{ordinfinitycharacter}،
$\cardeq{\beta}{\beta\ordplus1}$. اکنون بنا بر
\olref[cardinals][cardsasords]{lem:CardinalsBehaveRight}،
$\card{\beta}=\card{\beta\ordplus1}=\card{\alpha}$؛ پس
$\alpha\neq\card{\alpha}$.
\end{proof}

پیش‌تر، در~\olref[sfr][siz][enm-alt]{defn:enumerable}، اشاره کردیم که
می‌توان میان مجموعه‌های نامتناهیِ~!!{enumerable} و~!!{nonenumerable}
تمایز گذاشت. آن تعریف به‌طور طبیعی به نتیجهٔ زیر می‌انجامد:

\begin{prop}
$A$~!!{enumerable} است اگر و تنها اگر~$\card{A}\leq\omega$، و
$A$~!!{nonenumerable} است اگر و تنها اگر~$\omega<\card{A}$.
\end{prop}

\begin{proof}
بنا بر سه‌شقی، دو ادعا هم‌ارزند؛ پس کافی است اثبات کنیم~$A$
!!{enumerable} است اگر و تنها اگر~$\card{A}\leq\omega$. برای جهت
\emph{راست به چپ}: اگر~$\card{A}\leq\omega$، آنگاه بنا بر
\olref[cardinals][cardsasords]{lem:CardinalsBehaveRight} و
\olref{omegaisacardinal}، $\cardle{A}{\omega}$. برای جهت
\emph{چپ به راست}: فرض کنید~$A$~!!{enumerable} باشد؛ آنگاه بنا بر
\olref[sfr][siz][enm-alt]{defn:enumerable} سه حالت ممکن وجود دارد:
\begin{enumerate}
	\item اگر~$A=\emptyset$، آنگاه بنا بر~\olref{naturalsarecardinals}
	و~\olref[cardinals][cardsasords]{lem:CardinalsBehaveRight}،
	$\card{A}=0\in\omega$؛
	\item اگر~$\cardeq{n}{A}$، آنگاه بنا بر
	\olref{naturalsarecardinals} و
	\olref[cardinals][cardsasords]{lem:CardinalsBehaveRight}،
	$\card{A}=n\in\omega$؛
	\item اگر~$\cardeq{\omega}{A}$، آنگاه بنا بر
	\olref{omegaisacardinal}، $\card{A}=\omega$.
\end{enumerate}
پس در همهٔ حالت‌ها~$\card{A}\leq\omega$.
\end{proof}
\noindent
در واقع، $\omega$ جایگاهی ویژه دارد. هرچند اعداد ترتیبی شمارای بسیاری
وجود دارند:

\begin{cor}
$\omega$ تنها کاردینال نامتناهیِ~!!{enumerable} است.
\end{cor}

\begin{proof}
فرض کنید~$\cardfont{a}$ یک کاردینال نامتناهیِ~!!a{enumerable} باشد.
چون~$\cardfont{a}$ نامتناهی است، $\omega\leq\cardfont{a}$. چون
$\cardfont{a}$ کاردینالی~!!a{enumerable} است،
$\cardfont{a}=\card{\cardfont{a}}\leq\omega$. پس بنا بر سه‌شقی،
$\cardfont{a}=\omega$.
\end{proof}

البته بی‌نهایت کاردینال وجود دارد. پس شاید بپرسیم: \emph{چند کاردینال
وجود دارد؟} نتایج زیر نشان می‌دهند شاید بهتر باشد در این پرسش بازاندیشی
کنیم.

\begin{prop}\ollabel{unioncardinalscardinal}
اگر هر عضو~$X$ یک کاردینال باشد، آنگاه~$\bigcup X$ یک کاردینال است.
\end{prop}

\begin{proof}
به‌آسانی می‌توان بررسی کرد که~$\bigcup X$ یک عدد ترتیبی است. فرض کنید
$\alpha\in\bigcup X$ یک عدد ترتیبی باشد؛ آنگاه برای یک کاردینال
$\cardfont{b}$، $\alpha\in\cardfont{b}\in X$. چون~$\cardfont{b}$ یک
کاردینال است، $\cardless{\alpha}{\cardfont{b}}$. چون
$\cardfont{b}\subseteq\bigcup X$، داریم
$\cardle{\cardfont{b}}{\bigcup X}$؛ و بنابراین
$\cardneq{\alpha}{\bigcup X}$. با تعمیم، $\bigcup X$ یک کاردینال است.
\end{proof}

\begin{thm}\ollabel{lem:NoLargestCardinal}
بزرگ‌ترین کاردینال وجود ندارد.
\end{thm}

\begin{proof}
برای هر کاردینال~$\cardfont{a}$، قضیهٔ کانتور
(\olref[sfr][siz][car]{thm:cantor}) و
\olref[cardinals][cardsasords]{lem:CardinalsExist} نتیجه می‌دهند
$\cardfont{a}<\card{\Pow{\cardfont{a}}}$.
\end{proof}

\begin{thm}
مجموعهٔ همهٔ کاردینال‌ها وجود ندارد.
\end{thm}

\begin{proof}
برای برهان خلف فرض کنید
$C=\Setabs{\cardfont{a}}{\cardfont{a}\text{ یک کاردینال است}}$.
اکنون بنا بر~\olref{unioncardinalscardinal}، $\bigcup C$ یک کاردینال
است؛ پس بنا بر~\olref{lem:NoLargestCardinal} کاردینالی مانند
$\cardfont{b}>\bigcup C$ وجود دارد. بنا بر تعریف، $\cardfont{b}\in C$؛
پس~$\cardfont{b}\subseteq\bigcup{C}$ و در نتیجه
$\cardfont{b}\leq\bigcup C$، که تناقض است.
\end{proof}

شایسته است این نتیجه را هم با پارادوکس راسل و هم با پارادوکس
بورالی--فورتی مقایسه کنید.

\end{document}
